\documentclass[11pt]{article}
\usepackage{mathblatt}
% gesetzt von werkzeuge/setzer.py aus dem Lernweg (L4-1 · L4-2 · L4-3 · L4-4 · L4-5 · L4-6) und den Bankzeilen; Muster T6B.
% Präambel des Setzers (werkzeuge/setzer.py), wortgleich aus dem Hand-Satz T6B (bau/proben/2026-10-09/kathete-berechnen), dort aus K4W.
\makeatletter
\setlength{\mbpfbe}{0mm}% keine Punkte (Bauregel 6.4)
% eingerückt (Bauregel 4.7, 6.6): \gEin Einzug, \gSchrift eine Stufe kleiner
\newlength{\gEin}\setlength{\gEin}{0pt}
\let\gSchrift\relax
\renewcommand{\pfkreuzzeile}[1]{\par\noindent\hangindent3.2em\hangafter1\hspace*{1.6em}$\square$\ #1\par}
\newcommand{\gkreuz}[1]{$\square$\ #1\hspace{2em}}
% Bauregel 6.7: links, was man ansieht, rechts, was man tut; Spaltengrenze überall an derselben Stelle
\newsavebox{\gbox}\newlength{\gLinks}
\newcommand{\gzwei}[2]{\par\smallskip\noindent\sbox{\gbox}{#1}%
  \setlength{\gLinks}{\dimexpr0.5\textwidth-\gEin-\mbpfnr\relax}%
  \ifdim\wd\gbox>\dimexpr\gLinks-4mm\relax
    \usebox{\gbox}\par\smallskip\noindent\begin{minipage}{\linewidth}#2\end{minipage}%
  \else
    \begin{minipage}[t]{\gLinks}\vspace{0pt}\usebox{\gbox}\end{minipage}%
    \begin{minipage}[t]{\dimexpr\linewidth-\gLinks\relax}\vspace{0pt}\raggedright #2\end{minipage}%
  \fi\par}
\RenewDocumentEnvironment{pfaufg}{m m m}{%
  \begin{lrbox}{\mbaufgabebox}\begin{minipage}[t]{\linewidth}%
  \gSchrift\noindent\hspace*{\gEin}\mbpfrandmarke{#1}%
  \begin{minipage}[t]{\mbpfnr}\strut\textbf{#2}\end{minipage}%
  \begin{minipage}[t]{\dimexpr\linewidth-\gEin-\mbpfnr\relax}\raggedright\strut\ignorespaces}%
 {\end{minipage}%
  \end{minipage}\end{lrbox}%
  \setlength{\mbaufgabehoehe}{\dimexpr\ht\mbaufgabebox+\dp\mbaufgabebox\relax}%
  \par\addvspace{7pt}\Needspace*{\mbaufgabehoehe}%
  \noindent\usebox{\mbaufgabebox}\par\addvspace{7pt}}
\makeatother
% Rechenraum als Karo (Bauregel 6.9): n Rechenschritte = n mal 10 mm
\newcommand{\karo}[1]{\par\smallskip\noindent\begin{tikzpicture}
  \clip (0,0) rectangle ({\linewidth-1mm},{#1*10mm});
  \draw[black!16,line width=0.3pt,step=5mm] (0,0) grid ({\linewidth},{#1*10mm});
  \end{tikzpicture}\par}
\newcommand{\kk}[1]{$\square$\,#1\hspace{0.9em}}
\newcommand{\linie}[1]{\,{\color{black!45}\rule[-1pt]{#1}{0.4pt}}\,}
% Zeile am Ende jedes Durchgangs: Originale klein grau, darüber; dann Vorgänger/Nachfolger (Bauregel 3.4, 6.5)
\newcommand{\blattende}[2]{\par\vfill\noindent\begin{minipage}{\linewidth}{\scriptsize\color{mbgrau}Originale: #1\par}\smallskip
{\footnotesize #2\par}\end{minipage}}
\newcommand{\pkt}{\textperiodcentered{}}

\newcommand{\kennung}{WFW}
\newcommand{\grau}[1]{{\color{mbgrau}#1}}

\begin{document}
\pfheftstil
\fancyfoot[C]{\parbox[t]{\textwidth}{\mbfhbox\par\vspace{1.5mm}{\footnotesize\color{mbgrau}{\scriptsize \kennung}\hfill\thepage}}}
\pfheftkopf{Grundwert berechnen}{Klasse 7}
% L4-1: Anschluss an Einheit 3, rückwärts: Jetzt ist der graue Teil bekannt und der ganze Streifen gesucht; der Teil passt 2-, 4-, 5- oder 10-mal hinein, das Ganze ist größer als der Teil; in d) passt der Teil nicht glatt (8 % sind 6 €), ohne Streifen – Vervielfachen reicht nicht mehr
\begin{pfaufg}{}{1.}{}
Der graue Teil ist gegeben. Wie groß ist das Ganze ($100\,\%$)?\par\medskip\noindent
\begin{minipage}[b]{0.245\linewidth}\centering
\begin{tikzpicture}
\fill[black!28] (0,0) rectangle (1.800,0.6);
\draw[black!55,line width=0.4pt] (1.800,0) -- (1.800,0.6);
\draw[line width=0.7pt] (0,0) rectangle (3.6,0.6);
\draw[line width=0.4pt] (0,0.68) -- (0,0.82) -- (1.800,0.82) -- (1.800,0.68);
\node[above right,font=\small,inner xsep=0pt] at (0,0.80) {$50\,\%$: $18$ kg};
\node[below,font=\small] at (1.8,0) {ganz: ?};
\end{tikzpicture}\par\smallskip
\grau{{\small a)\ $50\,\%$ sind $18$ kg}\par\smallskip
$100\,\% = 2 \cdot 18$\par\smallskip $=$ $36$ kg}
\end{minipage}\hfill
\begin{minipage}[b]{0.245\linewidth}\centering
\begin{tikzpicture}
\fill[black!28] (0,0) rectangle (0.900,0.6);
\draw[black!55,line width=0.4pt] (0.900,0) -- (0.900,0.6);
\draw[black!55,line width=0.4pt] (1.800,0) -- (1.800,0.6);
\draw[black!55,line width=0.4pt] (2.700,0) -- (2.700,0.6);
\draw[line width=0.7pt] (0,0) rectangle (3.6,0.6);
\draw[line width=0.4pt] (0,0.68) -- (0,0.82) -- (0.900,0.82) -- (0.900,0.68);
\node[above right,font=\small,inner xsep=0pt] at (0,0.80) {$25\,\%$: $7\,€$};
\node[below,font=\small] at (1.8,0) {ganz: ?};
\end{tikzpicture}\par\smallskip
{\small b)\ $25\,\%$ sind $7\,€$}\par\smallskip
\linie{30mm}
\end{minipage}\hfill
\begin{minipage}[b]{0.245\linewidth}\centering
\begin{tikzpicture}
\fill[black!28] (0,0) rectangle (0.720,0.6);
\draw[black!55,line width=0.4pt] (0.720,0) -- (0.720,0.6);
\draw[black!55,line width=0.4pt] (1.440,0) -- (1.440,0.6);
\draw[black!55,line width=0.4pt] (2.160,0) -- (2.160,0.6);
\draw[black!55,line width=0.4pt] (2.880,0) -- (2.880,0.6);
\draw[line width=0.7pt] (0,0) rectangle (3.6,0.6);
\draw[line width=0.4pt] (0,0.68) -- (0,0.82) -- (0.720,0.82) -- (0.720,0.68);
\node[above right,font=\small,inner xsep=0pt] at (0,0.80) {$20\,\%$: $9$ m};
\node[below,font=\small] at (1.8,0) {ganz: ?};
\end{tikzpicture}\par\smallskip
{\small c)\ $20\,\%$ sind $9$ m}\par\smallskip
\linie{30mm}
\end{minipage}\hfill
\begin{minipage}[b]{0.245\linewidth}\centering
\begin{tikzpicture}
\path (0,-0.35) rectangle (3.6,0.95);
\end{tikzpicture}\par\smallskip
{\small d)\ $8\,\%$ sind $6\,€$}\par\smallskip
\linie{30mm}
\end{minipage}
\fusshilfe{\mbox{1:~b) 28\,€; c) 45 m; d) 75\,€}}
\end{pfaufg}

% L4-2: Erkennen ohne Rechnen: „25 % von 80 €“ und „25 % sind 80 €“ sehen gleich aus, verlangen aber Gegenteiliges; Kontrolle: Teil gesucht – kleiner, Ganzes gesucht – größer
\begin{pfaufg}{}{2.}{}
Kreuze die Rechnung an, die passt. Rechne nicht aus.\par\medskip
\grau{a)\ $25\,\%$ von $80\,€$. Wie viel Euro sind das?\par\hspace*{1.2em}$\boxtimes$\,$80 : 4$\hspace{0.9em}$\square$\,$80 \cdot 4$\hspace{0.9em}}\par\medskip
b)\ $25\,\%$ sind $80\,€$. Wie viel Euro ist das Ganze?\par\hspace*{1.2em}\kk{$80 : 4$}\kk{$80 \cdot 4$}\par\medskip
c)\ $10\,\%$ sind $35$ kg. Wie viel kg ist das Ganze?\par\hspace*{1.2em}\kk{$35 \cdot 10$}\kk{$35 : 10$}\kk{$0{,}1 \cdot 35$}\par\medskip
d)\ $10\,\%$ von $35$ kg. Wie viel kg sind das?\par\hspace*{1.2em}\kk{$35 \cdot 10$}\kk{$0{,}1 \cdot 35$}\kk{$10 : 35$}\par\medskip
e)\ Mia spart $20\,\%$. Das sind $9\,€$. Wie viel kostete die Ware vorher?\par\hspace*{1.2em}\kk{$0{,}2 \cdot 9$}\kk{$9 \cdot 5$}\kk{$9 : 5$}\par\medskip
f)\ Eine Jacke kostet $60\,€$. Es gibt $20\,\%$ Rabatt. Wie viel Euro spart man?\par\hspace*{1.2em}\kk{$60 \cdot 5$}\kk{$60 : 0{,}2$}\kk{$0{,}2 \cdot 60$}\par
\fusshilfe{\mbox{2:~b) $80 \cdot 4$; c) $35 \cdot 10$; d) $0{,}1 \cdot 35$; e) $9 \cdot 5$; f) $0{,}2 \cdot 60$}}
\end{pfaufg}

% L4-3: Jeder Satz geht über 1 %: Prozentwert durch den Prozentsatz, dann mal 100; zuerst mit den Zahlen aus 1c (20 % sind 9 m → wieder 45 m), dann glatte Sätze, zuletzt krumme mit dem Taschenrechner in einem Zug (W : p · 100); die Probe (p % vom Ergebnis gibt den Teil) fängt die Verwechslung
\begin{pfaufg}{}{3.}{}
Rechne über $1\,\%$: Prozentwert durch Prozentsatz, dann mal $100$. Bei krummen Zahlen nimm den Taschenrechner.\par\medskip
\grau{a)\ $20\,\%$ sind $9$ m\quad $1\,\% = 9\text{ m} : 20 = 0{,}45$ m, \ $100\,\% = 100 \cdot 0{,}45\text{ m} = 45$ m (wie in 1c)\quad Probe: $20\,\%$ von $45$ m $= 9$ m $\checkmark$}\par\medskip
\noindent\begin{minipage}[t]{0.48\linewidth}
b)\ $3\,\%$ sind $27$ kg
\karo{1}
\end{minipage}\hfill
\begin{minipage}[t]{0.48\linewidth}
c)\ $15\,\%$ sind $90\,€$
\karo{1}
\end{minipage}
\noindent\begin{minipage}[t]{0.48\linewidth}
d)\ $6\,\%$ sind $1{,}8$ l
\karo{1}
\end{minipage}\hfill
\begin{minipage}[t]{0.48\linewidth}
e)\ $36\,\%$ sind $63$ kg
\karo{1}
\end{minipage}
\noindent\begin{minipage}[t]{0.48\linewidth}
f)\ $2{,}5\,\%$ sind $3$ m
\karo{1}
\end{minipage}
\fusshilfe{\mbox{3:~b) 900 kg; c) 600\,€; d) 30 l; e) 175 kg; f) 120 m}}
\end{pfaufg}

% L4-4: In einer Sache stehen Prozentsatz und Prozentwert anders da (Akku-Anzeige mit Restzeit); das Ganze berechnen und damit entscheiden – einmal reicht es, einmal nicht
\begin{pfaufg}{}{4.}{}
Ben fährt $16$ Stunden mit dem Bus. Die Anzeigen zeigen, wie viel Akku noch da ist und wie lange er damit reicht. Reicht ein voll geladener Akku für die ganze Fahrt?\par\medskip\noindent
\begin{minipage}[b]{0.48\linewidth}\centering
\begin{tikzpicture}
\fill[black!28] (0,0) rectangle (0.840,1.0);
\draw[line width=0.8pt,rounded corners=1.5pt] (0,0) rectangle (2.4,1.0);
\fill (2.4,0.32) rectangle (2.52,0.68);
\node[font=\small\bfseries] at (1.2,0.5) {$35\,\%$};
\node[below,font=\small] at (1.2,-0.05) {noch $7$ Stunden};
\end{tikzpicture}\par\smallskip
{\small a)\ Kopfhörer}\par\smallskip
\linie{30mm}\par\medskip\linie{30mm}
\end{minipage}\hfill
\begin{minipage}[b]{0.48\linewidth}\centering
\begin{tikzpicture}
\fill[black!28] (0,0) rectangle (0.960,1.0);
\draw[line width=0.8pt,rounded corners=1.5pt] (0,0) rectangle (2.4,1.0);
\fill (2.4,0.32) rectangle (2.52,0.68);
\node[font=\small\bfseries] at (1.2,0.5) {$40\,\%$};
\node[below,font=\small] at (1.2,-0.05) {noch $6$ Stunden};
\end{tikzpicture}\par\smallskip
{\small b)\ Handy}\par\smallskip
\linie{30mm}\par\medskip\linie{30mm}
\end{minipage}
\fusshilfe{\mbox{4:~a) ja, 20 h; b) nein, 15 h}}
\end{pfaufg}

% L4-5: Gemischt mit Einheit 2 und 3: zu jeder Aufgabe erst hinschreiben, was gesucht ist (Prozentwert, Prozentsatz oder Grundwert), dann rechnen und die Probe mit der Gegenrechnung machen
\begin{pfaufg}{}{5.}{}
Schreib zuerst auf, was gesucht ist: Prozentwert, Prozentsatz oder Grundwert. Dann rechne und mach die Probe mit der Gegenrechnung.\par\medskip
\grau{a)\ $15$ von $60$ Losen gewinnen. Wie viel Prozent sind das?\quad gesucht: Prozentsatz; $15 : 60 = 0{,}25 = 25\,\%$\quad Probe: $25\,\%$ von $60 = 15$ $\checkmark$}\par\medskip
\noindent\begin{minipage}[t]{0.48\linewidth}
b)\ $30\,\%$ der $40$ Kinder einer Klasse fahren mit dem Bus. Wie viele Kinder sind das?
\karo{2}
\end{minipage}\hfill
\begin{minipage}[t]{0.48\linewidth}
c)\ $18$ Kinder fahren mit dem Rad. Das sind $60\,\%$ der Klasse. Wie viele Kinder hat die Klasse?
\karo{2}
\end{minipage}
\noindent\begin{minipage}[t]{0.48\linewidth}
d)\ Ein Buch hat $240$ Seiten. Lea hat $84$ Seiten gelesen. Wie viel Prozent sind das?
\karo{2}
\end{minipage}\hfill
\begin{minipage}[t]{0.48\linewidth}
e)\ Zum Sommerfest kommen $135$ Gäste. Das sind $90\,\%$ der Eingeladenen. Wie viele waren eingeladen?
\karo{2}
\end{minipage}
\noindent\begin{minipage}[t]{0.48\linewidth}
f)\ Ein Rad kostet $480\,€$. Es gibt $15\,\%$ Rabatt. Wie viel Euro spart man?
\karo{2}
\end{minipage}
\fusshilfe{\mbox{5:~b) 12 Kinder; c) 30 Kinder; d) 35\,\%; e) 150 Gäste; f) 72\,€}}
\end{pfaufg}

% L4-6: Zielaufgabe: In einer Rabatt-Tabelle wechselt die Frage von Zeile zu Zeile; der alte Preis ist immer das Ganze, der Rabatt in Euro der Teil; aus Rabatt in Euro und in Prozent den alten Preis berechnen; beim alten Preis die Probe mit der Gegenrechnung
\begin{pfaufg}{nach P10 ’25}{6.}{}
Im Schlussverkauf hängen Preisschilder. In jeder Zeile fehlen zwei Angaben. Ergänze die Tabelle. Mach beim alten Preis die Probe mit der Gegenrechnung.\par\medskip\noindent {\renewcommand{\arraystretch}{1.3}\begin{tabular}{|l|c|c|c|c|}\hline
Ware & alter Preis & Rabatt in \% & Rabatt in € & neuer Preis\\\hline
Schal & $40\,€$ & $25\,\%$ &  & \\\hline
Jacke & $90\,€$ &  & $18\,€$ & \\\hline
Schuhe &  & $20\,\%$ & $14\,€$ & \\\hline
Rucksack &  & $15\,\%$ & $9\,€$ & \\\hline\end{tabular}}\par\smallskip
\noindent\begin{minipage}[t]{0.48\linewidth}
a)\ Schal: Rabatt in Euro und neuer Preis
\karo{2}
\end{minipage}\hfill
\begin{minipage}[t]{0.48\linewidth}
b)\ Jacke: Rabatt in Prozent und neuer Preis
\karo{2}
\end{minipage}
\noindent\begin{minipage}[t]{0.48\linewidth}
c)\ Schuhe: alter Preis und neuer Preis
\karo{2}
\end{minipage}\hfill
\begin{minipage}[t]{0.48\linewidth}
d)\ Rucksack: alter Preis und neuer Preis
\karo{2}
\end{minipage}
\fusshilfe{\mbox{6:~a) 10\,€; 30\,€; b) 20\,\%; 72\,€; c) 70\,€; 56\,€; d) 60\,€; 51\,€}}
\end{pfaufg}

\par\vfill\noindent{\footnotesize Vorher: Prozentwert berechnen\quad\textbullet\quad Weiter: Prozentuale Veränderung\par}
\end{document}
